
\documentclass[aspectratio=169,11pt]{beamer}

% ------------------------------------------------------------
% Presentación Beamer
% Módulo 5: Stress Testing Actuarial y Demográfico
% Curso MACROPOL: Stress Testing en Sistemas de Pensiones
% ------------------------------------------------------------

\usepackage[utf8]{inputenc}
\usepackage[T1]{fontenc}
\usepackage[spanish,es-nodecimaldot]{babel}
\usepackage{lmodern}
\usepackage{amsmath, amssymb}
\usepackage{booktabs}
\usepackage{array}
\usepackage{multicol}
\usepackage{tikz}
\usepackage{xcolor}
\usepackage{listings}
\usepackage{hyperref}

\usetheme{Madrid}
\usecolortheme{seahorse}
\setbeamertemplate{navigation symbols}{}
\setbeamertemplate{footline}[frame number]

% ------------------------------------------------------------
% Colores
% ------------------------------------------------------------
\definecolor{macroblue}{RGB}{0,55,105}
\definecolor{macrogray}{RGB}{70,70,70}
\definecolor{macrolight}{RGB}{232,240,248}
\definecolor{macrored}{RGB}{150,35,35}
\definecolor{macrogreen}{RGB}{35,115,70}

\setbeamercolor{title}{fg=white,bg=macroblue}
\setbeamercolor{frametitle}{fg=macroblue}
\setbeamercolor{structure}{fg=macroblue}
\setbeamercolor{block title}{fg=white,bg=macroblue}
\setbeamercolor{block body}{bg=macrolight,fg=black}
\setbeamercolor{alerted text}{fg=macrored}

% ------------------------------------------------------------
% Configuración de código R
% ------------------------------------------------------------
\lstdefinestyle{Rstyle}{
  language=R,
  basicstyle=\ttfamily\scriptsize,
  keywordstyle=\color{macroblue}\bfseries,
  commentstyle=\color{macrogreen},
  stringstyle=\color{macrored},
  numbers=left,
  numberstyle=\tiny\color{macrogray},
  stepnumber=1,
  numbersep=6pt,
  backgroundcolor=\color{macrolight},
  showspaces=false,
  showstringspaces=false,
  breaklines=true,
  frame=single,
  rulecolor=\color{macroblue}
}

% ------------------------------------------------------------
% Comandos útiles
% ------------------------------------------------------------
\newcommand{\baseline}{\textbf{\color{macrogreen}Baseline}}
\newcommand{\adverso}{\textbf{\color{macrored}Adverso}}
\newcommand{\supuestos}{\textbf{Supuestos}}
\newcommand{\resultados}{\textbf{Resultados}}
\newcommand{\narrativa}{\textbf{Narrativa económica}}
\newcommand{\modelo}{\textbf{Modelo cuantitativo}}

\title[Stress Testing Actuarial]{Módulo 5\\Stress Testing Actuarial y Demográfico}
\subtitle{Curso: Stress Testing en Sistemas de Pensiones de Capitalización Individual}
\author{MACROPOL: Economía Aplicada y Políticas Públicas}
\date{}

\begin{document}

% ------------------------------------------------------------
\begin{frame}
  \titlepage
  \vspace{-0.3cm}
  \begin{center}
    \small Impacto de longevidad, densidad de cotización y cambios demográficos sobre la suficiencia previsional.
  \end{center}
\end{frame}

% ------------------------------------------------------------
\begin{frame}{Contexto del módulo}
\begin{block}{Ubicación dentro del curso}
\textbf{Módulo 5: Stress Testing Actuarial y Demográfico}
\end{block}

\begin{itemize}
  \item El Módulo 4 cuantifica shocks de portafolio.
  \item El Módulo 5 traduce esos shocks a \textbf{pensiones futuras}.
  \item El Módulo 6 integra resultados financieros, actuariales y regulatorios.
\end{itemize}

\vspace{0.2cm}
\begin{alertblock}{Pregunta central}
¿Cómo cambian las tasas de reemplazo cuando la población vive más, cotiza menos y enfrenta retornos más bajos?
\end{alertblock}
\end{frame}

% ------------------------------------------------------------
\begin{frame}{Objetivos de aprendizaje}
Al finalizar el módulo, los participantes podrán:

\begin{enumerate}
  \item Identificar riesgos actuariales relevantes en cuentas individuales.
  \item Construir escenarios demográficos baseline y adversos.
  \item Medir impactos sobre fondo acumulado, pensión y tasa de reemplazo.
  \item Implementar simulaciones actuariales simples en R.
  \item Interpretar resultados para supervisión, regulación y política pública.
\end{enumerate}
\end{frame}

% ------------------------------------------------------------
\begin{frame}{Agenda de la sesión}
\begin{enumerate}
  \item Fundamentos del riesgo actuarial.
  \item Riesgo de longevidad.
  \item Densidad de cotización e informalidad.
  \item Tasa de reemplazo como indicador de suficiencia.
  \item Simulación actuarial y Monte Carlo en R.
  \item Escenarios demográficos integrados.
  \item Taller aplicado y discusión regulatoria.
\end{enumerate}
\end{frame}

% ------------------------------------------------------------
\section{Contexto}

\begin{frame}{Contexto: pensiones bajo capitalización individual}
\begin{center}
\begin{tikzpicture}[node distance=1.9cm, every node/.style={font=\small}]
\node[draw, rounded corners, fill=macrolight, minimum width=2.2cm, minimum height=0.8cm] (c) {Cotizaciones};
\node[draw, rounded corners, fill=macrolight, right of=c, xshift=1.4cm, minimum width=2.4cm, minimum height=0.8cm] (f) {Fondo acumulado};
\node[draw, rounded corners, fill=macrolight, right of=f, xshift=1.6cm, minimum width=2.3cm, minimum height=0.8cm] (r) {Retiro};
\node[draw, rounded corners, fill=macrolight, right of=r, xshift=1.2cm, minimum width=2.4cm, minimum height=0.8cm] (p) {Pensión};
\draw[->, thick, macroblue] (c) -- (f);
\draw[->, thick, macroblue] (f) -- (r);
\draw[->, thick, macroblue] (r) -- (p);
\end{tikzpicture}
\end{center}

\vspace{0.3cm}
\begin{block}{Intuición económica}
En un sistema de cuentas individuales, la suficiencia previsional depende de tres bloques: cuánto se aporta, cuánto rinde el ahorro y durante cuántos años debe financiarse la pensión.
\end{block}
\end{frame}

% ------------------------------------------------------------
\begin{frame}{Riesgo actuarial: definición operativa}
\begin{block}{Definición}
Riesgo de que los supuestos demográficos, laborales y actuariales usados para proyectar beneficios futuros no se materialicen.
\end{block}

\begin{columns}
\column{0.48\textwidth}
\textbf{Fuentes principales}
\begin{itemize}
  \item Mayor longevidad.
  \item Menor densidad de cotización.
  \item Jubilación anticipada.
  \item Crecimiento salarial bajo.
  \item Brechas de género y empleo.
\end{itemize}

\column{0.48\textwidth}
\textbf{Resultados afectados}
\begin{itemize}
  \item Fondo acumulado.
  \item Pensión inicial.
  \item Tasa de reemplazo.
  \item Probabilidad de insuficiencia.
  \item Costo fiscal de mitigación.
\end{itemize}
\end{columns}
\end{frame}

% ------------------------------------------------------------
\begin{frame}{Riesgo financiero vs. riesgo actuarial}
\begin{table}
\centering
\begin{tabular}{p{0.42\textwidth}p{0.42\textwidth}}
\toprule
\textbf{Riesgo financiero} & \textbf{Riesgo actuarial} \\
\midrule
Horizonte corto/mediano & Horizonte largo \\
Precios de activos & Mortalidad y supervivencia \\
Volatilidad de retornos & Años esperados en retiro \\
Tasas de interés & Factores de conversión actuarial \\
Pérdida de valor del portafolio & Menor pensión para un mismo saldo \\
\bottomrule
\end{tabular}
\end{table}

\begin{alertblock}{Mensaje clave}
Un shock financiero puede ser transitorio; un aumento permanente de longevidad cambia estructuralmente la relación entre ahorro acumulado y pensión mensual.
\end{alertblock}
\end{frame}

% ------------------------------------------------------------
\section{Modelo / Herramienta}

\begin{frame}{Modelo básico de acumulación}
\begin{block}{Cuenta individual}
\[
F_T = \sum_{t=1}^{T} c \cdot d_t \cdot W_t \cdot \prod_{j=t}^{T}(1+r_j)
\]
\end{block}

\begin{itemize}
  \item $F_T$: fondo acumulado al retiro.
  \item $c$: tasa de contribución.
  \item $d_t$: densidad de cotización.
  \item $W_t$: salario cotizable.
  \item $r_j$: rentabilidad real del portafolio.
\end{itemize}

\vspace{0.2cm}
\textbf{Intuición:} pequeñas diferencias persistentes en densidad y rentabilidad se acumulan exponencialmente durante la vida laboral.
\end{frame}

% ------------------------------------------------------------
\begin{frame}{Modelo básico de desacumulación}
\begin{block}{Conversión del saldo en pensión}
\[
P = \frac{F_T}{a_x}
\]
\[
a_x = \sum_{k=1}^{\omega-x} \frac{{}_{k}p_x}{(1+i)^k}
\]
\end{block}

\begin{itemize}
  \item $P$: pensión anual inicial.
  \item $a_x$: factor actuarial de anualidad a la edad $x$.
  \item ${}_{k}p_x$: probabilidad de sobrevivir $k$ años desde edad $x$.
  \item $i$: tasa técnica de descuento.
\end{itemize}

\begin{alertblock}{Canal de estrés}
Mayor supervivencia aumenta $a_x$ y reduce $P$ para un mismo fondo acumulado.
\end{alertblock}
\end{frame}

% ------------------------------------------------------------
\begin{frame}{Tasa de reemplazo}
\begin{block}{Indicador de suficiencia}
\[
TR = \frac{P}{W_T}
\]
\end{block}

\begin{columns}
\column{0.5\textwidth}
\textbf{Interpretación}
\begin{itemize}
  \item Alta: preserva consumo en retiro.
  \item Media: requiere ahorro complementario.
  \item Baja: riesgo de vulnerabilidad social.
\end{itemize}

\column{0.45\textwidth}
\begin{table}
\centering
\begin{tabular}{cc}
\toprule
TR & Lectura \\
\midrule
$>$ 70\% & Adecuada \\
50--70\% & Intermedia \\
$<$ 50\% & Vulnerable \\
\bottomrule
\end{tabular}
\end{table}
\end{columns}
\end{frame}

% ------------------------------------------------------------
\begin{frame}{Baseline vs. adverso}
\begin{columns}
\column{0.48\textwidth}
\begin{block}{\baseline}
\begin{itemize}
  \item Rentabilidad real: 5\%.
  \item Densidad de cotización: 80\%.
  \item Edad de retiro: 65.
  \item Longevidad según tabla vigente.
  \item Crecimiento salarial real: 1.5\%.
\end{itemize}
\end{block}

\column{0.48\textwidth}
\begin{alertblock}{\adverso}
\begin{itemize}
  \item Rentabilidad real: 2\%.
  \item Densidad de cotización: 60\%.
  \item Edad efectiva de retiro: 62.
  \item Longevidad: +5 años.
  \item Crecimiento salarial real: 0\%.
\end{itemize}
\end{alertblock}
\end{columns}
\end{frame}

% ------------------------------------------------------------
\begin{frame}{Narrativa económica vs. modelo cuantitativo}
\begin{table}
\centering
\begin{tabular}{p{0.44\textwidth}p{0.44\textwidth}}
\toprule
\narrativa & \modelo \\
\midrule
Envejecimiento más rápido que el previsto. & Shock a probabilidades de supervivencia ${}_{k}p_x$. \\
Aumento de informalidad laboral. & Caída de $d_t$ en la etapa de acumulación. \\
Retornos persistentemente bajos. & Reducción de $r_t$ en todas las trayectorias. \\
Retiro anticipado por deterioro laboral. & Menor $T$ y mayor horizonte de pago. \\
\bottomrule
\end{tabular}
\end{table}

\begin{block}{Regla pedagógica}
La narrativa define plausibilidad; el modelo traduce esa historia en pérdidas previsionales cuantificables.
\end{block}
\end{frame}

% ------------------------------------------------------------
\begin{frame}{Supuestos vs. resultados}
\begin{columns}
\column{0.48\textwidth}
\textbf{\supuestos}
\begin{itemize}
  \item Rentabilidad.
  \item Densidad.
  \item Edad de retiro.
  \item Mortalidad.
  \item Salario.
  \item Tasa técnica.
\end{itemize}

\column{0.48\textwidth}
\textbf{\resultados}
\begin{itemize}
  \item Fondo acumulado.
  \item Pensión anual.
  \item Tasa de reemplazo.
  \item Percentil 5 de TR.
  \item Probabilidad de TR $<$ umbral.
  \item Brecha contra baseline.
\end{itemize}
\end{columns}
\end{frame}

% ------------------------------------------------------------
\section{Riesgo de longevidad}

\begin{frame}{Riesgo de longevidad}
\begin{block}{Definición}
Riesgo de que los afiliados vivan más que lo incorporado en las tablas actuariales usadas para calcular beneficios o rentas vitalicias.
\end{block}

\begin{itemize}
  \item Aumenta el número esperado de pagos.
  \item Reduce la pensión inicial si el saldo acumulado no cambia.
  \item Eleva el costo de garantías mínimas o subsidios complementarios.
  \item Puede tensionar aseguradoras que ofrecen rentas vitalicias.
\end{itemize}
\end{frame}

% ------------------------------------------------------------
\begin{frame}{Shock de longevidad}
\begin{table}
\centering
\begin{tabular}{lccc}
\toprule
Escenario & Esperanza de vida al retiro & Horizonte de pago & Efecto esperado \\
\midrule
Baseline & 20 años & Normal & Referencia \\
Adverso & 25 años & +5 años & Menor pensión \\
Severo & 28 años & +8 años & Mayor brecha \\
\bottomrule
\end{tabular}
\end{table}

\begin{alertblock}{Intuición}
Con el mismo fondo acumulado, más años esperados de retiro implican una pensión anual menor.
\end{alertblock}
\end{frame}

% ------------------------------------------------------------
\begin{frame}{Ejemplo numérico: anualidad simple}
Suponga un saldo al retiro de RD\$10,000,000 y tasa técnica real de 3\%.

\[
P = \frac{F_T}{\sum_{k=1}^{n}(1+i)^{-k}}
\]

\begin{table}
\centering
\begin{tabular}{lccc}
\toprule
Escenario & $n$ años & Factor anualidad & Pensión anual \\
\midrule
Baseline & 20 & 14.88 & 672,157 \\
Adverso & 25 & 17.41 & 574,279 \\
Severo & 28 & 18.76 & 532,932 \\
\bottomrule
\end{tabular}
\end{table}

\textbf{Lectura:} el shock de longevidad reduce la pensión anual aunque el saldo no cambie.
\end{frame}

% ------------------------------------------------------------
\section{Densidad de cotización}

\begin{frame}{Densidad de cotización}
\begin{block}{Definición}
Porcentaje de la vida laboral en que el trabajador efectivamente realiza aportes al sistema.
\end{block}

\begin{itemize}
  \item 100\%: cotización continua.
  \item 80\%: interrupciones moderadas.
  \item 60\%: informalidad o desempleo frecuente.
  \item 50\% o menos: trayectoria altamente vulnerable.
\end{itemize}

\begin{alertblock}{Canal previsional}
Una menor densidad reduce el capital acumulado y limita el efecto de capitalización compuesta.
\end{alertblock}
\end{frame}

% ------------------------------------------------------------
\begin{frame}{Stress test de densidad}
\begin{table}
\centering
\begin{tabular}{lccc}
\toprule
Escenario & Densidad & Narrativa & Resultado esperado \\
\midrule
Baseline & 80\% & Formalidad estable & TR de referencia \\
Adverso & 65\% & Mayor informalidad & Menor acumulación \\
Severo & 50\% & Intermitencia laboral & Riesgo de insuficiencia \\
\bottomrule
\end{tabular}
\end{table}

\begin{block}{Decisión regulatoria}
La densidad de cotización permite evaluar si las reglas de aportación producen pensiones suficientes para trayectorias laborales reales, no solo para carreras completas.
\end{block}
\end{frame}

% ------------------------------------------------------------
\section{Implementación}

\begin{frame}[fragile]{Implementación en R: parámetros}
\begin{lstlisting}[style=Rstyle]
# Parametros de base
edad_inicio  <- 25
edad_retiro  <- 65
anios        <- edad_retiro - edad_inicio

salario0     <- 1000
g_salario    <- 0.015
tasa_aporte  <- 0.10
retorno      <- 0.05
densidad     <- 0.80

tasa_tecnica <- 0.03
vida_retiro  <- 20
\end{lstlisting}
\end{frame}

% ------------------------------------------------------------
\begin{frame}[fragile]{Implementación en R: acumulación}
\begin{lstlisting}[style=Rstyle]
# Funcion de acumulacion previsional
acumular_fondo <- function(salario0, g, aporte, densidad,
                           retorno, anios) {
  fondo <- 0
  salario <- salario0

  for (t in 1:anios) {
    contribucion <- salario * aporte * densidad
    fondo <- (fondo + contribucion) * (1 + retorno)
    salario <- salario * (1 + g)
  }

  return(list(fondo = fondo, salario_final = salario))
}
\end{lstlisting}
\end{frame}

% ------------------------------------------------------------
\begin{frame}[fragile]{Implementación en R: pensión y tasa de reemplazo}
\begin{lstlisting}[style=Rstyle]
# Factor de anualidad simple
factor_anualidad <- function(i, n) {
  sum(1 / (1 + i)^(1:n))
}

# Calculo de pension y tasa de reemplazo
res <- acumular_fondo(salario0, g_salario, tasa_aporte,
                      densidad, retorno, anios)

ax <- factor_anualidad(tasa_tecnica, vida_retiro)
pension_anual <- res$fondo / ax
tr <- pension_anual / (12 * res$salario_final)

tr
\end{lstlisting}
\end{frame}

% ------------------------------------------------------------
\begin{frame}[fragile]{Implementación en R: escenarios}
\begin{lstlisting}[style=Rstyle]
escenarios <- data.frame(
  nombre = c("Baseline", "Adverso", "Severo"),
  retorno = c(0.05, 0.02, 0.00),
  densidad = c(0.80, 0.65, 0.50),
  vida_retiro = c(20, 25, 28)
)

calcular_escenario <- function(r, d, n) {
  res <- acumular_fondo(salario0, g_salario, tasa_aporte,
                        d, r, anios)
  pension <- res$fondo / factor_anualidad(tasa_tecnica, n)
  tr <- pension / (12 * res$salario_final)
  c(fondo = res$fondo, pension_anual = pension, TR = tr)
}

resultados <- t(mapply(calcular_escenario,
                       escenarios$retorno,
                       escenarios$densidad,
                       escenarios$vida_retiro))
cbind(escenarios, round(resultados, 3))
\end{lstlisting}
\end{frame}

% ------------------------------------------------------------
\begin{frame}[fragile]{Monte Carlo actuarial}
\begin{lstlisting}[style=Rstyle]
set.seed(123)
N <- 10000

retornos  <- rnorm(N, mean = 0.04, sd = 0.08)
densidades <- pmin(pmax(rnorm(N, mean = 0.75, sd = 0.12), 0.30), 1)
vida_ret  <- round(pmax(rnorm(N, mean = 22, sd = 3), 15))

simular_TR <- function(r, d, n) {
  res <- acumular_fondo(salario0, g_salario, tasa_aporte,
                        d, r, anios)
  pension <- res$fondo / factor_anualidad(tasa_tecnica, n)
  pension / (12 * res$salario_final)
}

TR_sim <- mapply(simular_TR, retornos, densidades, vida_ret)

summary(TR_sim)
quantile(TR_sim, probs = c(0.05, 0.50, 0.95))
mean(TR_sim < 0.50)
\end{lstlisting}
\end{frame}

% ------------------------------------------------------------
\begin{frame}{Indicadores de salida}
\begin{table}
\centering
\begin{tabular}{p{0.36\textwidth}p{0.52\textwidth}}
\toprule
Indicador & Uso supervisor \\
\midrule
TR promedio & Suficiencia esperada del sistema \\
Percentil 5 de TR & Cola de riesgo previsional \\
Probabilidad de TR $<$ 50\% & Riesgo de insuficiencia masiva \\
Brecha contra baseline & Magnitud del shock adverso \\
Fondo acumulado promedio & Sensibilidad a aportes y retornos \\
Pensión inicial & Impacto directo sobre afiliados \\
\bottomrule
\end{tabular}
\end{table}
\end{frame}

% ------------------------------------------------------------
\section{Ejercicio}

\begin{frame}{Ejercicio aplicado: diseño de escenarios}
\begin{block}{Trabajo en equipos}
Construir un stress test actuarial para una cohorte de trabajadores que inicia cotización a los 25 años y se retira a los 65.
\end{block}

\begin{table}
\centering
\begin{tabular}{lccc}
\toprule
Variable & Baseline & Adverso & Severo \\
\midrule
Rentabilidad real & 5\% & 2\% & 0\% \\
Densidad & 80\% & 65\% & 50\% \\
Crecimiento salarial & 1.5\% & 0.5\% & 0\% \\
Vida en retiro & 20 años & 25 años & 28 años \\
\bottomrule
\end{tabular}
\end{table}
\end{frame}

% ------------------------------------------------------------
\begin{frame}{Ejercicio aplicado: tareas}
Cada equipo debe calcular:

\begin{enumerate}
  \item Fondo acumulado al retiro.
  \item Pensión anual inicial.
  \item Tasa de reemplazo.
  \item Pérdida de TR frente al baseline.
  \item Probabilidad de TR inferior a 50\%.
\end{enumerate}

\vspace{0.2cm}
\begin{alertblock}{Entrega}
Una tabla de resultados y una recomendación regulatoria de no más de cinco líneas.
\end{alertblock}
\end{frame}

% ------------------------------------------------------------
\begin{frame}{Preguntas guía para discusión}
\begin{enumerate}
  \item ¿Qué variable explica más la caída de la tasa de reemplazo?
  \item ¿El riesgo de longevidad o la baja densidad genera mayor vulnerabilidad?
  \item ¿Qué indicador debería monitorear el regulador trimestralmente?
  \item ¿Qué medidas de política reducen la cola de riesgo previsional?
  \item ¿Cómo cambia la conclusión si se reduce la tasa técnica?
\end{enumerate}
\end{frame}

% ------------------------------------------------------------
\section{Interpretación}

\begin{frame}{Interpretación de resultados}
\begin{block}{Lectura técnica}
El stress test no busca predecir el futuro exacto, sino cuantificar la vulnerabilidad del sistema bajo condiciones plausibles y severas.
\end{block}

\begin{itemize}
  \item \textbf{Baseline:} punto de comparación, no meta de política.
  \item \textbf{Adverso:} combinación coherente de shocks demográficos y financieros.
  \item \textbf{Resultados:} distribución de tasas de reemplazo, no solo promedio.
\end{itemize}
\end{frame}

% ------------------------------------------------------------
\begin{frame}{Implicaciones regulatorias}
\begin{itemize}
  \item Revisar tablas de mortalidad y supuestos de longevidad.
  \item Evaluar suficiencia de tasas de contribución.
  \item Diseñar alertas tempranas por cohortes vulnerables.
  \item Revisar edad efectiva de retiro y densidad real de cotización.
  \item Integrar resultados actuariales con stress testing de portafolio.
\end{itemize}

\begin{alertblock}{Decisión de política pública}
Cuando la cola de la distribución de TR cae por debajo de umbrales mínimos, el regulador debe evaluar reformas paramétricas, incentivos a formalidad o mecanismos de pensión mínima.
\end{alertblock}
\end{frame}



% ------------------------------------------------------------
\begin{frame}{Conexión con decisiones de política}
\begin{table}
\centering
\begin{tabular}{p{0.42\textwidth}p{0.48\textwidth}}
\toprule
Hallazgo del stress test & Decisión posible \\
\midrule
TR baja por longevidad & Actualizar tablas y factores actuariales \\
TR baja por densidad & Incentivos a formalidad y cotización continua \\
Alta dispersión por retornos & Revisar ciclo de vida e inversión por edad \\
Riesgo alto en mujeres & Ajustar política de equidad previsional \\
Probabilidad elevada de TR $<$ 50\% & Evaluar pensión mínima o ahorro complementario \\
\bottomrule
\end{tabular}
\end{table}
\end{frame}

% ------------------------------------------------------------
\begin{frame}{Mensajes clave}
\begin{enumerate}
  \item La longevidad transforma el riesgo previsional en un problema estructural.
  \item La densidad de cotización puede ser tan importante como la rentabilidad.
  \item La tasa de reemplazo debe ser el indicador principal de suficiencia.
  \item El análisis debe distinguir supuestos, escenarios y resultados.
  \item La supervisión moderna requiere integrar portafolio, demografía y política pública.
\end{enumerate}
\end{frame}

% ------------------------------------------------------------
\begin{frame}{Referencias de apoyo}
\small
\begin{itemize}
  \item Impavido, G. ``Conducting Stress Tests of Defined Benefit Pension Plans'', en \textit{A Guide to IMF Stress Testing Methods and Models}, FMI.
  \item IOPS. \textit{Stress Testing and Scenario Analysis of Pension Plans}, Working Paper No. 19.
  \item OECD. \textit{Pensions at a Glance}: indicadores de tasas de reemplazo.
  \item Banco Mundial. Estudios sobre tasas de reemplazo, cuentas individuales y fase de desacumulación.
  \item Material base del curso MACROPOL: \textit{Stress Testing en Sistemas de Pensiones de Capitalización Individual}.
\end{itemize}
\end{frame}

% ------------------------------------------------------------
\begin{frame}
\begin{center}
{\Huge Gracias}

\vspace{0.5cm}

{\Large Discusión y preguntas}

\vspace{0.5cm}

\textbf{Módulo 5: Stress Testing Actuarial y Demográfico}
\end{center}
\end{frame}

\end{document}
